%% Supplemental material for the TCC manuscript (range-only scheme).
%% Proofs adapted from paper/jisa/supplementary.tex (S2-S4 there), with the
%% keyword, retrieval, flexible and role parts removed.
\documentclass[10pt,journal,compsoc]{IEEEtran}
\usepackage{amsmath,amssymb,amsthm}
\usepackage[nocompress]{cite}
\newtheorem{theorem}{Theorem}
\newtheorem{lemma}{Lemma}
\newtheorem{definition}{Definition}
\newtheorem*{theorem*}{Theorem}
\renewcommand{\thesection}{S\arabic{section}}
% Shared by the LaTeX build (manuscript.tex) and the Word build (build_docx.py).
\newcommand{\negl}{\mathsf{negl}}
\newcommand{\PPT}{\mathrm{PPT}}
\newcommand{\Real}{\mathbf{Real}}
\newcommand{\Ideal}{\mathbf{Ideal}}
\newcommand{\Leak}{\mathcal{L}}
\newcommand{\Sim}{\mathcal{S}}
\newcommand{\Adv}{\mathcal{A}}
\newcommand{\cind}{\approx_c}
\newcommand{\Gen}{\mathsf{Gen}}
\newcommand{\Eval}{\mathsf{Eval}}
\newcommand{\Enc}{\mathsf{Enc}}
\newcommand{\Dec}{\mathsf{Dec}}
\newcommand{\advPRF}{\mathrm{Adv}^{\mathrm{prf}}}
\newcommand{\advDPF}{\mathrm{Adv}^{\mathrm{dpf}}}
\newcommand{\advCPA}{\mathrm{Adv}^{\mathrm{cpa}}}
\newcommand{\bplus}{\boxplus}
\newcommand{\bdot}{\boxdot}


\begin{document}
\title{Supplemental Material for ``One-Round Verifiable Range Aggregation over
Encrypted Data Replicated across Non-Colluding Clouds''}
\author{Md~Mazharul~Islam, Mubasshir~Ahmed, and Md~Redowan~Zaman~Anik}
\maketitle

\noindent Section, equation and theorem numbers without an ``S'' refer to the
main paper. Notation follows Table~2 of the main paper; $u_x$ denotes the unit
vector at cell $x$.

% ---------------------------------------------------------------------------
\section{Proof of Theorem~1 (Correctness)}

\begin{theorem*}[Theorem 1, restated]
Suppose $w_f[i]\,m_i\ge0$ for every record and $\sum_iw_f[i]\,m_i<p$. Then a
query on box $B$ returns $\sum_{i:\,c(i)\in B}w_f[i]\,m_i$, except with the
probability that decryption fails, and honest replies pass the check.
\end{theorem*}

\begin{IEEEproof}
For each corner $\sigma$ the analyst generates keys for the point function
$f_{\kappa_\sigma(B),\beta_\sigma}$, with $\beta_\sigma=(-1)^{|\sigma|}\bmod p$
for a corner inside the cube and $\beta_\sigma=0$ at a random cell for a dummy.
By DPF correctness, the evaluations of all $n$ clouds sum to the point
function, so
\[
e:=\sum_{j}e_j=\sum_{\sigma}\beta_\sigma\,u_{\kappa_\sigma(B)},
\]
where dummies contribute the zero vector whatever cell they name. Cloud $j$
places its evaluation of the $\beta$-keys in row~0 and of the
$\delta_s\beta$-keys in row~1 of the same plaintexts, and multiplies them into
the encrypted table, which holds $T_f$ in both rows. By Lemma~1 the sum of the
replies decrypts in row~0 to
\[
\sum_x e[x]\,T_f[x]=\sum_{\sigma\text{ real}}(-1)^{|\sigma|}\,T_f[\kappa_\sigma(B)]\pmod p .
\]
The corners outside the cube are exactly those that would contribute $0$ in
Eq.~(2), since a summed-area table is zero below the origin. Eq.~(2) therefore
turns the right-hand side into the box total $Y=\sum_{c\in B}C_f[c]$ modulo
$p$. Because every contribution is non-negative, $0\le Y\le\sum_iw_f[i]\,m_i<p$,
so reduction modulo $p$ returns $Y$ exactly, although single signed terms are
negative.

\emph{Segments.} Several templates share a ciphertext, each in its own segment
of each row. Rotations move values only within a row, and the rotate-and-sum
covers one segment, so the slot of template $r$ receives only template $r$'s
products, and the two rows do not mix. Compressing the reply to one modulus
tower is a public map that preserves the decrypted value.

\emph{Check.} In row~1 the keys share $\delta_s$ times the same point
functions, so the same argument gives $\delta_sY$. Honest replies therefore
pass the test $\hat y=\delta_sy$ with certainty. The only failure term is the
decryption-failure probability of the leveled scheme under the noise growth
discussed in the proof of Lemma~1.
\end{IEEEproof}

% ---------------------------------------------------------------------------
\section{Proof of Theorem~2 (Privacy)}

\emph{Leakage.} The setup leakage is
$\Leak_{\mathrm{Setup}}=(n,p,\ell,\{(s_1,\dots,s_d)\}_{\text{cubes}},
\text{number of templates and ciphertexts})$. The query leakage of query $q$ is
$\Leak_{\mathrm{Query}}(q)=(\text{cube addressed},\ \text{check keys present or
not})$. It fixes the number of keys a cloud receives, $2^d$ or $2^{d+1}$.

\begin{theorem*}[Theorem 2, restated]
If $(\Gen,\Eval)$ is an $(n{-}1)$-private DPF and BFV is IND-CPA secure given
its evaluation keys, then for every efficient adversary corrupting at most
$n-1$ clouds there is an efficient simulator such that
\[
\big|\Pr[\Real=1]-\Pr[\Ideal=1]\big|\le Q\,2^{d_{\max}+1}\,\advDPF+m\,\advCPA ,
\]
where $Q$ is the number of queries, $d_{\max}$ the largest cube dimension and
$m$ the number of ciphertexts the coalition stores.
\end{theorem*}

\begin{IEEEproof}
Take a coalition $C$ of exactly $n-1$ clouds; a smaller coalition sees a
subset of the same view. Its view consists of the stored ciphertexts, the
public and evaluation keys, and for each query its members' keys and the
replies they compute. Let $H_0$ be the real experiment.

\emph{$H_1$: simulated keys.} Replace, one key generation at a time, the keys
of the coalition by the output of the DPF simulator for the right domain $[M]$.
$(n{-}1)$-privacy hides the output $\beta$ as well as the point, so this
covers real corners, dummy corners and the keys carrying $\delta_s\beta$ alike.
A query to a $d$-dimensional cube makes $2^d$ key generations, and $2^d$ more
when it carries the check, so the hybrid costs at most
$Q\,2^{d_{\max}+1}\advDPF$. After it, the keys depend only on the domain and
on how many keys the query carries, which the query leakage fixes. In
particular they no longer depend on the box, on which corners were dummies, or
on $\delta_s$.

\emph{$H_2$: encryptions of zero.} Replace each of the $m$ stored ciphertexts
by an encryption of $0$ under $pk$. The coalition holds only $pk$ and the
evaluation keys, so by IND-CPA security this costs at most $m\,\advCPA$. As is
standard for BFV, we assume IND-CPA security holds given the rotation keys
(circular security).

\emph{Simulator.} In $H_2$ every stored value is an encryption of $0$ whose
number and shape are given by $\Leak_{\mathrm{Setup}}$, and every key is
simulated from $\Leak_{\mathrm{Query}}$. A reply is a deterministic function
of a cloud's keys, the addressed table and the public keys, so the simulator
recomputes the replies of the coalition itself. It thus produces $H_2$ from
the leakage alone, which is the ideal experiment. Summing the two hybrid costs
proves the bound.
\end{IEEEproof}

\emph{Why sharing the selection matters.} Suppose instead that the analyst
recovered the set of matching cells and asked the clouds to sum over it. That
set is a function of the query and the data, and no simulator given only the
leakage above could produce it. The proof works because each cloud contracts
its own share of the selection, which is pseudorandom on its own.

\emph{Deviating clouds.} Clouds never interact and receive nothing but keys,
so a coalition that deviates sees the same view as one that does not. Its
deviation gains it no information about queries; Theorem~3 makes any
deviation that changes an answer evident. Not covered are output delivery and
the rejection events themselves, discussed after Theorem~3 in the main paper.

% ---------------------------------------------------------------------------
\section{Proof of Theorem~3 (Detection of wrong answers)}

\begin{theorem*}[Theorem 3, restated]
Let up to $n-1$ clouds return arbitrary replies computed from their views,
with at least one cloud honest, $\delta_s$ known only to the analyst, and no
cloud colluding with the analyst. The analyst accepts an output $y'\neq y$
with probability at most $1/(p-1)+2^d\advDPF$ per query, and the bounds add
over $Q$ queries.
\end{theorem*}

\begin{IEEEproof}
Fix one query and one template. Let $V$ be the coalition's view: both key sets
of its members, the stored ciphertexts, the public and evaluation keys, and
everything it saw in earlier queries. Its replies, summed over $C$, decrypt to
some $(r_0,r_1)$ in the two rows, a function of $V$, the data and the
coalition's coins. The decryption key is independent of $\delta_s$. Let $y_C$
and $\hat y_C$ be what the coalition would have contributed honestly, also
functions of $V$ and the data. The honest clouds contribute $y-y_C$ and
$\delta_sy-\hat y_C$. The analyst therefore decrypts
\[
y'=y-y_C+r_0,\qquad \hat y'=\delta_sy-\hat y_C+r_1,
\]
and accepts if and only if $\hat y'=\delta_sy'$, that is,
\begin{equation}
\label{eq:accept}
r_1-\hat y_C=\delta_s\,(y'-y).
\end{equation}

Replace the coalition's $2^d$ keys that carry $\delta_s\beta_\sigma$ by keys
for the same points and an independent scalar $\delta'_s$. A reduction that
knows the points, the signs, the data and $\delta_s$ can evaluate the event
``$y'\neq y$ and \eqref{eq:accept} holds'' from the coalition's keys and
outputs. By $(n{-}1)$-privacy, which hides the output as well as the point,
the probability of the event changes by at most $2^d\advDPF$.

After the change, $V$ is independent of $\delta_s$. Conditioned on $V$, the
data and the coins, every quantity in \eqref{eq:accept} except $\delta_s$ is
fixed. Since $y'-y\neq0$, the right-hand side is a bijective function of
$\delta_s$, so at most one of the $p-1$ equally likely values of $\delta_s$
satisfies \eqref{eq:accept}. This gives the bound for one query. Each query
draws a fresh $\delta_s$, so what the coalition learns from earlier queries,
including whether they were accepted, concerns independent scalars, and a
union bound over $Q$ queries completes the proof.
\end{IEEEproof}

The proof places no restriction on how a cloud deviates. It may change its
share, rescale the queried cell, swap or mix the two rows, replay an old reply
or return any ciphertext. The bound does not guarantee an answer, since a
deviating cloud can always force a rejection. Whether the analyst rejects can
depend on the data, for example on whether a cell the cloud chose to change is
empty. An analyst that reveals its rejections can therefore leak up to one bit
about the data per query, though nothing about the query. Retrying silently
through another set of clouds closes this channel.

\emph{Several secrets.} With $r$ independent secrets
$\delta_s^{(1)},\dots,\delta_s^{(r)}$, each with its own key set and slot
segment, the analyst accepts only if all $r$ relations hold. Repeating the
argument for each secret, with the other secrets fixed, gives
$(p-1)^{-r}+r\,2^d\advDPF$. Our implementation uses $r=1$.

% ---------------------------------------------------------------------------
\section{Why the Data Are Encrypted Rather Than Secret-Shared}

Under requirements (R1)--(R3) of the main paper, one might deal the data as an
additive sharing $T=\sum_kS_k$ and let cloud $j$ answer
$\langle e_j,S_j\rangle$. This is wrong: $\langle e,T\rangle=\sum_{j,k}\langle
e_j,S_k\rangle$, and the cross terms $j\neq k$ are held by no single cloud.
For $n=2$ the error is $-\langle e_1,T\rangle+\langle e_1-e_2,S_1\rangle$,
which is uniform because $S_1$ is uniform and independent of the keys whenever
$e_1\neq e_2$.

The cross terms are not an artefact of linear answers. Take any
non-interactive scheme in which cloud~1 holds an additive share $S_1$ of the
data, the client forms queries without looking at the data, and the scheme is
correct for every tuple of shares. Given a database $X$ held by one server, a
client can run the scheme with $S_1:=X$ and the other shares chosen by itself.
It plays clouds $2,\dots,n$ locally, sends only the first query to the server,
and subtracts the other contributions. Privacy against cloud~1 hides the query
from the server, so this is single-server private information retrieval, which
implies oblivious transfer~\cite{dicrescenzo2000single} and is usually built
from homomorphic encryption~\cite{kushilevitz1997replication}. Within this
class, sharing the data does not remove the public-key primitive; it moves it
into a retrieval per share.

The argument does not cover all designs. A lower threshold escapes it: with
$(t{+}1,n)$ shares and $t<n/2$, servers can multiply shared values locally,
which is the route of Obscure~\cite{gupta2020obscure} at the price of an
honest majority and $2l+1$ servers for a predicate of length $l$. Interaction
between servers or preprocessing also escape it, and (R2) excludes the first.
Encryption under one owner key has operational advantages as well. A shared
store with threshold $n-1$ is lost when one share is lost, whereas copies
survive the loss of $n-1$ clouds. Queries continue over any two surviving
clouds. The owner can add or replace a cloud by copying ciphertexts instead of
re-dealing the whole store.

% ---------------------------------------------------------------------------
\section{Check Layouts with Stored Keys}

The evaluation compares the per-query check with two layouts that store keys.

\emph{One stored key.} Row~1 holds $\delta T_f$ for a secret $\delta$, and
the analyst accepts only if $\hat y=\delta y$. A cloud that adds a constant to
one row, corrupts one row or returns an unrelated ciphertext is caught except
with probability $1/p$. A cloud that applies the same linear change to both
rows passes. Contracting against $k\,e_j$ instead of $e_j$, for example,
shifts $y$ and $\hat y$ by amounts in the ratio $\delta$.

\emph{Per-cell stored keys.} Row~1 holds $\delta_cT_f[c]$ with a key per
cell, and the reply carries one segment per corner so that the analyst can
check each corner against its own key. This needs $2^d$ copies of the table.
A change $e_j+\Delta$ then passes only if
$\sum_{c}\Delta[c](\delta_c-\delta_{\kappa})T_f[c]=0$ at each checked corner
$\kappa$, a non-zero linear form in independent keys whenever $\Delta$ touches
a non-empty cell other than $\kappa$. A change at the queried corner itself
leaves the ratio intact and passes, which a cloud that knows or guesses the
corner can exploit. Table~3 of the main paper confirms both failures; the
per-query check catches them at the cost of a second key set and no extra
homomorphic product.

\bibliographystyle{IEEEtran}
\bibliography{references}
\end{document}
